%% =====================================================================
%%  B1-T-08-02 -- what B1 contributes to "The Engagement Decision"
%%  -------------------------------------------------------------------
%%  LAYER A: APPEND-ONLY. Written once, then left alone. Material found
%%  only here sits in the `unique` slot and is protected by rule R10.
%% =====================================================================
%% @kind: source
%% @id: B1-T-08-02
%% @book: B1
%% @topic: T-08-02
%% @locator: Tactic 5, pp. 59--70; Tactic 10, pp. 117--124
%% @status: distilled
%% @unique: yes
%% @integrated: yes
%% @updated: 2026-09-19

\dossier{B1}{T-08-02}{\bookshort{B1}}{Tactic 5, pp. 59--70; Tactic 10, pp. 117--124}

\begin{distilled}
  \item Not retaliating is argued on cost grounds: an ego injury creates a persistent enemy and the friction piles problems on you, emotionally and financially.
  \item A rule of thumb is given for which arguments to contest: let sleeping dogs lie where they do not affect your vital interests.
  \item Most arguments are said to be winnable, which makes the decision to have one a decision rather than a necessity.
\end{distilled}

\begin{unique}
  \item The vital-interests filter as an explicit criterion for choosing which engagements to enter, with everything outside it conceded by default.
\end{unique}
